% ===== 原文 pp.44–45 =====
% [原文 p.44]

它定义一个代数环面 $\mathbf{T}^{N-k}_{Q,t}\subset(\C^\times)^N$，与 $(\C^\times)^{N-k}$ 同构。于是剩下的是扭超势

\begin{equation}
	\tW = \mu\sum_{i=1}^{N}\ee^{-y_i}. \tag{5.28}
\end{equation}

也就是说，环面 sigma 模型的镜像是靶空间为 $\mathbf{T}^{N-k}_{Q,t}$、超势为 $\tW=\mu\sum_{i=1}^{N}\ee^{-y_i}$ 的 Landau–Ginzburg 模型。

例如，$\CP^1$ 模型的镜像是靶空间为 $\mathbf{T}^1_{(1,1),t}$ 的 Landau–Ginzburg 模型，后者通过 $(y_1,y_2)=(y,t-y)$ 同构于 $\C^\times\ni y$，超势为 $\tW=\mu\ee^{-y}+\mu\ee^{-t+y}$。用 RG 不变的变量 $\ty=y-t/2$ 与 RG 不变的标度参数 $\Lambda=\ee^{-t/2}\mu$ 来写，它就是

\begin{equation}
	\tW = \Lambda\left(\ee^{-\ty}+\ee^{\ty}\right). \tag{5.29}
\end{equation}

% [Tutor Note: (5.29)–(5.30) 的 RG 不变改写可直接核对：取 $\ty_i=y_i-t/N$、$\Lambda=\mu\ee^{-t/N}$，则 $\mu\ee^{-y_i}=\Lambda\ee^{-\ty_i}$；又由约束 $\sum_{i=1}^{N}y_i=t$ 得 $\sum_{i=1}^{N-1}\ty_i=t/N-y_N$，故 (5.30) 的末项 $\mu\ee^{-t+y_1+\cdots+y_{N-1}}=\mu\ee^{-y_N}=\Lambda\ee^{\ty_1+\cdots+\ty_{N-1}}$。(5.29) 正是 $N=2$ 的特例，此时 $\Lambda=\mu\ee^{-t/2}$，与正文所给的一致。]

这就是 §3.3 中所主张的镜像对称。类似地，$\CP^{N-1}$ 模型的镜像是靶空间为 $(\C^\times)^{N-1}$、超势为 $\tW=\mu(\ee^{-y_1}+\cdots+\ee^{-y_{N-1}}+\ee^{-t+y_1+\cdots+y_{N-1}})$ 的 Landau–Ginzburg 模型，或者用 RG 不变的变量与参数写为

\begin{equation}
	\tW = \Lambda\left(\ee^{-\ty_1}+\cdots+\ee^{-\ty_{N-1}}+\ee^{\ty_1+\cdots+\ty_{N-1}}\right), \tag{5.30}
\end{equation}

\subsection{镜像对称的推导二——有超势的理论}

让我们考虑模型 $\mathrm{T}^{U(1)}_{N,d}$。暂且设超势为零，$W=0$。那么我们知道，它等价于扭超势为

\begin{equation}
	\tW = \mu\left(\ee^{-y_1}+\cdots+\ee^{-y_N}+\ee^{-y_P}\right). \tag{5.31}
\end{equation}

的系统，其中 $y_1,\dots,y_N,y_P$ 是取值于 $\C/2\pi\ii\Z$ 的 A-手征变量，服从约束 $y_1+\cdots+y_N-dy_P=t$，即

\begin{equation}
	\ee^{-dy_P} = \ee^{t}\ee^{-y_1}\cdots\ee^{-y_N}. \tag{5.32}
\end{equation}

现在让我们把原理论中的超势 $W=pf(x)$ 找回来。这不会改变扭超势 (5.31)，因为 B-手征参数无法进入其中。然而，它确实以相当剧烈的方式改变了 Kähler 势，这就迫使我们更换场变量。有一种办法可以找到合适的变量，用到所谓的 D-膜中心荷。按照它，在 $\mathfrak{M}_W$ 中的某个特殊点上，好的变量是 $\C$ 值变量 $\tx_1,\dots,\tx_N$，它们与 $y$ 的关系为

\begin{equation}
	\ee^{-y_1} = \tx_1^d,\quad\cdots\quad \ee^{-y_N} = \tx_N^d,\quad \ee^{-y_P} = \ee^{t/d}\tx_1\cdots\tx_N. \tag{5.33}
\end{equation}

% [Tutor Note: (5.33) 与约束 (5.32) 的自洽性一行即可验证：由 $\ee^{-y_P}=\ee^{t/d}\tx_1\cdots\tx_N$ 得 $\ee^{-dy_P}=\ee^{t}\prod_i\tx_i^{d}=\ee^{t}\prod_i\ee^{-y_i}$，正是 (5.32)；把 (5.33) 代入 (5.31) 就直接得到下一页的 (5.34)。另注意 $\tx\mapsto y$ 不是单射：保持 (5.34) 不变的变换为 $\tx_i\mapsto\omega_i\tx_i$，其中 $\omega_i^{d}=1$ 且 $\prod_i\omega_i=1$，共 $d^{N-1}$ 个元素，这正是 (5.34) 之后所说的轨形群 $(\Z_d)^{N-1}$。]

% [原文 p.45]

注意这些变量确实满足约束 (5.32)。用新变量写出的扭超势 (5.31) 为

\begin{equation}
	\tW = \mu\left(\tx_1^d+\cdots+\tx_N^d+\ee^{t/d}\tx_1\cdots\tx_N\right). \tag{5.34}
\end{equation}

注意从 $\tx$ 到 $y$ 的映射不是一一的。为修正这一点，我们需要把 $(\tx_1,\dots,\tx_N)$ 与 $(\omega_1\tx_1,\dots,\omega_N\tx_N)$ 等同起来，其中 $\omega_1^d=\cdots=\omega_N^d=\omega_1\cdots\omega_N=1$。也就是说，模型 $\mathrm{T}^{U(1)}_{N,d}$ 的镜像是 Landau–Ginzburg 轨形 $(\C^N/(\Z_d)^{N-1},(5.34))$。

\subsection{例子}

让我们用一些例子来讨论镜像理论的物理。

\para{$\CP^{N-1}$ 模型}

% [Tutor Note: 原文此处小标题 "The $\CP^{N-1}$ Model" 之后没有任何正文——本讲义是不完整版（作者在首页自述 "A full version will be available soon"），故此处照原文只排小标题，未补写任何内容。$\CP^{N-1}$ 模型的镜像超势已在 (5.30) 显式给出，可与本标题对照阅读。]

\para{模型 $\mathrm{H}_n$}

模型 $\mathrm{H}_n$ 的镜像是 $\C/2\pi\ii\Z$ 变量 $y_1,\dots,y_4$ 的 Landau–Ginzburg 模型，这些变量服从 $y_1+y_2-ny_4=t^1$ 与 $y_3+y_4=t^2$，超势为 $\tW=\mu\sum_{i=1}^{4}\ee^{-y_i}$。解出这些约束，比如取 $y_2=t^1+nt^2-y_1-ny_3$ 与 $y_4=t^2-y_3$，并记 $x_1:=\ee^{-y_1}$、$x_2:=\ee^{-y_3}$，我们就看到它是靶空间为 $\{(x_1,x_2)\}=(\C^\times)^2$、超势为

\begin{equation}
	\tW = \mu\left(x_1+\frac{q_1q_2^n}{x_1x_2^n}+x_2+\frac{q_2}{x_2}\right), \tag{5.35}
\end{equation}

的模型，其中 $q_a:=\ee^{-t_a}$。正如 §4.4 所讨论的，模型 $\mathrm{H}_n$ 被等同于靶空间为 $\mathbb{F}_n$（$n=0,1,2$ 的情形）或 $\WP^2_{1,1,n}$（$n\ge2$ 的情形）的 sigma 模型。因此，上述 LG 模型可以看作这些 sigma 模型的镜像。真空方程 $\dd\tW=0$ 给出

\begin{equation}
	x_1-\frac{q_1q_2^n}{x_1x_2^n}=0,\qquad -n\frac{q_1q_2^n}{x_1x_2^n}+x_2-\frac{q_2}{x_2}=0. \tag{5.36}
\end{equation}

在 $\mu x_1=\sigma_1$ 与 $\mu x_2=\sigma_2$ 之下它与 (4.35) 一致；这并不奇怪，因为 $x$-$\sigma$ 关系正是由 $\partial_{y_1}\tW=\partial_{y_3}\tW=0$ 得出的，其中 $\tW$ 是 (5.22) 中同时含有 $\sigma$ 与 $y$ 的那个。

% [Tutor Note: (5.36) 可由 (5.35) 验证：因 $\dd x_i/\dd y_i=-x_i$，有 $\partial_{y_1}\tW=-x_1\partial_{x_1}\tW$、$\partial_{y_3}\tW=-x_2\partial_{x_2}\tW$。$\partial_{y_1}$ 条件给出所印第一式；$\partial_{y_3}$ 条件给出 $n q_1q_2^n/(x_1x_2^n)-x_2+q_2/x_2=0$，即所印第二式（同解）；若改用 $\partial_{x_2}\tW=0$ 并代入第一式 $x_1^2x_2^n=q_1q_2^n$，则得等价的 $x_2^2-nx_1x_2-q_2=0$。另：原文 "constanits" 应为 "constraints"（下文解出的正是约束方程中的 $y_2,y_4$），"no surprize" 应为 "no surprise"；原文同段先用上标 $t^1,t^2$、后写下标 $q_a:=\ee^{-t_a}$，记号不统一，译文照原文保留。]

\para{$n=0$}

模型 $\mathrm{H}_0$ 被等同于靶空间为 $\mathbb{F}_0=\CP^1\times\CP^1$ 的 sigma 模型。镜像超势 (5.35) 确实就是两份 $\CP^1$ 的镜像超势之和。

\para{$n=1$}

模型 $\mathrm{H}_1$ 被等同于靶空间为 $\mathbb{F}_1$（即 $\CP^2$ 的一点爆破）的 sigma 模型。我们回忆，对于具有 $q_1=\Lambda/\mu$ 与 $q_2=q(\Lambda/\mu)^2$ 的理论，在 $q\ll1$
